% Generated draft: requires mathematical and source review.
\documentclass[11pt]{article}
\usepackage[margin=25mm]{geometry}
\usepackage{fontspec}
\setmainfont{Latin Modern Roman}
\usepackage{amsmath,amssymb,mathtools,unicode-math}
\setmathfont{Latin Modern Math}
\usepackage{enumitem,needspace,xurl,hyperref}
\setlength{\parindent}{0pt}
\setlength{\parskip}{5pt}
\setlength{\emergencystretch}{3em}
\newcommand{\N}{\mathbb{N}}
\newcommand{\Z}{\mathbb{Z}}
\newcommand{\Q}{\mathbb{Q}}
\newcommand{\R}{\mathbb{R}}
\newcommand{\C}{\mathbb{C}}
\newcommand{\F}{\mathbb{F}}
\newcommand{\A}{\mathbb{A}}
\newcommand{\PP}{\mathbb{P}}
\newcommand{\E}{\mathbb{E}}
\newcommand{\eps}{\varepsilon}
\newcommand{\dd}{\,\mathrm{d}}
\begin{document}
\section*{UnsolvedMath ID: 2303006}
\textbf{Source problem code:} catalogue entry
\subsection{Definitions and mathematical statement}
The source records a planar result attributed to Wolf: if an entire-plane function \(u\) in the intended class has an angular majorant \(f\),
\[
u(re^{i\theta})\le f(\theta)\qquad(0<r<\infty),
\qquad
\int_0^{2\pi}\log^+f(\theta)\,\dd\theta<\infty,
\]
then \(u\) is bounded above and hence constant. For a nonnegative majorant, \(\log^+t=\log\max(1,t)\). The excerpt does not state the class of \(u\), measurability/pointwise conventions for \(f\), or the complete Wolf reference.

The requested three-dimensional analogue would replace directions \(e^{i\theta}\) by \(\omega\in S^2\), with a directional bound
\[
u(r\omega)\le F(\omega)\qquad(0<r<\infty),
\]
and seek an appropriate condition on \(F\), relative to surface measure \(\sigma\), giving the intended boundedness or constancy conclusion. A direct candidate condition would be \(\int_{S^2}\log^+F\,\dd\sigma<\infty\), but the question asks for an analogue rather than asserting this particular condition sufficient. No unprovided exponent or integrability criterion is invented.

A crucial missing datum is whether \(u\) is harmonic (\(\Delta u=0\)) or subharmonic (upper semicontinuous with \(\Delta u\ge0\) in the distributional sense), or belongs to another class. The inference from boundedness above to constancy, which is valid for entire-plane subharmonic functions, does not have the same scope for subharmonic functions on \(\R^3\). Thus a precise three-dimensional theorem must specify both the class and whether its target is boundedness, constancy, or a modified conclusion. The supplied OPEN-TRIAGE summary reports the question open in the 2018 edition and no recent resolution located.
\subsection{Short English statement}
Find a three-dimensional counterpart of a planar theorem controlling a function by an angular majorant with integrable positive logarithm. The missing function class and the intended boundedness-versus-constancy conclusion are essential.
\subsection{Sources}
\begin{itemize}
\item Local UnsolvedMath registry, permanent ID \(2303006\), exact source code \texttt{catalogue entry}, ``Research Problems in Function Theory --- Problem 3.6.''
\item Hayman, \emph{Research Problems in Function Theory} (2018), Problem 3.6; \url{https://arxiv.org/abs/1809.07200}. Source-TeX extraction; recorded access 2026-07-29.
\item The source attributes the planar theorem to Wolf without a complete reference. The supplied OPEN-TRIAGE summary gives no independent current-status citation.
\end{itemize}
\end{document}
